\chapter*{符号及数学约定}
\addcontentsline{toc}{chapter}{符号与数学约定}

除非局部另行声明，本书以大写粗斜体表示矩阵及其批量表示，以小写粗斜体表示向量，以普通斜体表示标量和索引。矩阵元素 $W_{ij}$ 与向量分量 $x_i$ 使用普通斜体；矩阵 $\mat{W}$、向量 $\vect{x}$ 及其梯度采用相应的粗斜体。参数集合 $\vect{\theta}$ 作为展平向量时使用粗斜体，标量参数依局部定义使用普通斜体。形状中的轴顺序具有语义，交换轴必须同时改变后续运算。同一字母在不同学科中可能有不同含义，章内符号表与紧邻公式的定义优先。批量计算通常采用行向量约定；低秩适配等局部推导使用列向量时会明确说明。

\begin{longtable}{p{.18\linewidth}p{.73\linewidth}}
\toprule 符号 & 约定含义\\\midrule\endhead
$B,T,S$ & 批量大小、目标或当前序列长度、源序列长度；有效长度可以因样本而异。\\
$d,d_k,d_v$ & 模型隐藏维数、单头查询与键维数、单头值维数。\\
$H,H_q,H_{kv}$ & 注意力头数、查询头数、键值头数；视觉章节中的 $H,W$ 局部表示图像高宽。\\
$L,d_f$ & 网络层数、前馈网络中间维数；训练损失通常写为 $\mathcal L$，业务损失等局部记号另行定义。\\
$V,\mathcal V$ & 词表大小、词表集合；注意力中 $\mat{V}$ 表示值矩阵时，词表大小另加下标。\\
$\vect{\theta},\eta$ & 模型参数、学习率；采样温度通常用 $\tau$。\\
$x_{<t},x_{\le t}$ & 位置 $t$ 之前、不晚于位置 $t$ 的序列前缀。\\
$\Real^{m\times n}$ & 由 $m$ 行 $n$ 列实数矩阵构成的空间。\\
$\mat{X}^{\mathrm{T}},\odot$ & 转置、逐元素乘法；普通并置表示矩阵乘法或标量乘法。\\
$\mathbf1[\cdot]$ & 条件成立时取一、否则取零的指示函数；$\mathbf1$ 也可表示同维全一向量，依上下文区分。\\
$\mathbb E,\operatorname{Var}$ & 期望与方差；所有统计结论均依赖所声明的采样分布。\\
$\|\vect{x}\|_2,\|\mat{W}\|_F$ & 向量欧氏范数、矩阵 Frobenius 范数；矩阵 $\|\mat{W}\|_2$ 表示谱范数。\\
$\nabla_{\vect{\theta}},\dif x$ & 对参数向量的梯度、积分或微分中的微分记号。\\
$\pi(a\mid s),V^\pi,Q^\pi$ & 强化学习中的策略、状态价值与动作价值；与注意力矩阵 $\mat{Q},\mat{V}$ 分属不同语境。\\
$\gamma,\lambda$ & 强化学习中常表示折扣与迹衰减；排队中的 $\lambda$ 表示到达率，依章内定义。\\
$\beta_t,\alpha_t,\bar\alpha_t$ & 离散扩散中的噪声方差、单步信号保留量与累计保留量。\\
$t,\tau$ & 时间、位置或采样温度的局部记号；离散步与连续路径时间不能直接互换。\\
$q_{\vect{\phi}}(z\mid x),p_{\vect{\theta}}(x)$ & 变分近似后验与生成分布；下标指出可学习参数。\\
$\widehat x,\widehat p$ & 估计、预测或有限样本统计量，具体性质由构造方式决定。\\
$\varnothing$ & 空集合或协议定义的空条件；空条件并不总是零向量。\\
$\operatorname{stopgrad}(x)$ & 前向保留数值、反向阻断该路径梯度的运算。\\
\bottomrule
\end{longtable}

概率质量用于离散取值，概率密度用于连续变量。离散分布满足 $\sum_xp(x)=1$；密度满足 $\int f(x)\dif x=1$，单点密度值不是单点概率。Logit 指归一化之前的实数分数，Softmax 指沿指定轴进行的归一化指数变换。

圆周率 $\constpi$、自然常数 $\conste$ 与虚数单位 $\consti$ 使用正体。表示策略、概率或置换的 $\pi$，以及表示变量或索引的 $e,i$ 使用斜体。

本书使用自然对数。Kullback--Leibler 散度（KL 散度）定义为
\begin{equation*}
 D_{\mathrm{KL}}(p\|q)=\sum_xp(x)\log\frac{p(x)}{q(x)}.
\end{equation*}
约定 $p(x)=0$ 的项为零；若某处 $p(x)>0$ 而 $q(x)=0$，散度为正无穷。该量一般不对称，也不满足距离所需的全部性质。连续分布使用相应密度积分。

公式中的 $O(\cdot)$ 描述计算量或存储量随问题规模增大的渐近增长阶。实际运行时间取决于硬件性能、数据搬运、并行方式和实现开销，需在指定环境中测量。资源公式明确区分参数数量、元素数量、字节数与时间；GB 采用十进制，GiB 采用二进制。

形状说明描述广播后的逻辑尺寸，具体存储由实现决定。概率公式中的求和或积分范围以局部定义为准；条件分布需要说明条件事件和支持集。边界值导致零分母、退化分布或未定义指标时，需单独处理。
